\begin{titlepage}\centering
{\Large Carl Friedrich Gauss\par}\bigskip
{\LARGE\bfseries Werke\par}\bigskip
{\Large Band III\par}\bigskip
{\large Ein neuer Beweis des Satzes, dass jedes Polynom in einer Variablen in reelle Faktoren ersten oder zweiten Grades zerlegt werden kann\par}\bigskip
{\large Vollständige Arbeiten von 1799, 1815, 1816 und 1828\par}\bigskip
Moderne deutsche Übersetzung\par\bigskip
\raggedright
Es folgen die historischen Titelblätter und die vollständige Dissertation von 1799, Artikel 1–24, mit den zehn historischen Fußnoten von Gauss und allen vier Figuren. Es folgt der vollständige zweite Beweis, Artikel 1–20 mit beiden historischen Fußnoten und seinem Titelblatt. Er wurde 1815 vorgelegt und 1816 veröffentlicht. Es folgt der vollständige dritte Beweis, vorgelegt am 30. Januar 1816 und im selben Jahr veröffentlicht, mit Titelblatt, Einleitung, Artikeln 1–4 und historischer Fußnote. Es folgt die vollständige Arbeit Beweis eines algebraischen Lehrsatzes von 1828 über positive und negative Wurzeln und Vorzeichenwechsel und Vorzeichenfolgen der Koeffizienten, auch bei fehlenden Gliedern. Spätere Texte des dritten Bandes sind in diesem Leseband nicht enthalten. Historische Verweise behalten die Nummerierung der Vorlage bei.
\par\medskip Historische Vorlage: Carl Friedrich Gauss, \emph{Werke}, dritter Band, Göttingen, 1866; Dissertation erstmals 1799 bei C. G. Fleckeisen in Helmstedt erschienen. Die ursprünglich deutschen Titelblätter und die Arbeit von 1828 sind in modernem Deutsch wiedergegeben; ihre historischen Originale werden gesondert bewahrt.
\par\medskip Mit OpenAI GPT-6.1 Sol KI-übersetzt ins Deutsche; KI-erzeugter Satz.
\end{titlepage}
